\documentclass[10pt,a4paper]{amsart}
\usepackage[margin=19mm]{geometry}
\usepackage{amsmath,amssymb,mathtools}
\usepackage[scr=rsfs,cal=euler]{mathalfa}
\usepackage{xfrac}
\usepackage[unicode,hidelinks,bookmarks=false]{hyperref}
\newcommand{\ie}{i.e.\ }
\newcommand{\cf}{cf.\ }
\newcommand{\resp}{respectively\ }
\newcommand{\Subsection}{\subsection}
\providecommand{\nobreakdash}{\nobreak\hbox{-}}
\setlength{\emergencystretch}{2em}
\multlinegap0pt
\usepackage{bm}
\usepackage{bbm}
\usepackage{dsfont}
\usepackage{xspace}
\usepackage{cancel}
\usepackage{mathtools}
\usepackage{amsthm}
\makeatletter
\@ifundefined{theorem}{\newtheorem{theorem}{Theorem}}{}
\makeatother
\newcommand{\CZ}{\mathbb{C}\mathfrak{Z}}
\newcommand{\mI}{\mathbb{I}}
\newcommand{\zdu}{\mathfrak{D}\mspace{1mu}}
\title[Spectral Properties of the Zeon Combinatorial Laplacian]{Spectral Properties of the Zeon Combinatorial Laplacian}
\author{G. Stacey Staples}
\date{Source: arXiv:2502.05678v1 (CC BY 4.0)}
\begin{document}
\maketitle

\noindent\emph{Source and licence.} Spectral Properties of the Zeon Combinatorial Laplacian. Version arXiv:2502.05678v1 (CC BY 4.0), distributed under the Creative Commons Attribution 4.0 International licence, as recorded in the arXiv licence metadata for this exact version and shown on the abstract page https://arxiv.org/abs/2502.05678v1. Full rights notice: licenses/arxiv-2502.05678-CC-BY-4.0.txt.\par\smallskip

\noindent\textit{Editorial scope.} Counted unit: the Theorem whose source label is \texttt{laplacian eigenvalue} in the article (source0.tex lines 471--476, source label \texttt{laplacian eigenvalue}), delivered together with the article's own complete original proof of it (source0.tex lines 478--711, one \texttt{proof} environment of 234 lines). No mathematics is written, repaired or re-derived: statement and proof are the article's own text line for line. Required context retained uncounted: none. Every definition, display and label of those lines is retained unchanged. Omitted from this package and not redistributed: 1--470, 477--477, 712--1314. Nothing inside a retained excerpt is omitted or altered: every retained body is delivered exactly as the article's lines are written, so no omission receipt of any kind accompanies this package. Citations: the retained excerpts carry no citation command at all, so no bibliography entry of the article is transcribed and no reference is redistributed. Reference adaptation: none was needed and none was made; every reference inside the delivered unit resolves inside it, so no unresolved reference or citation is produced. Printed slips kept literal: no printed word, symbol, hypothesis or number of the counted statement or of its proof was altered. Layout and class: the article's own class is \texttt{article} with packages including authblk, amsmath, amssymb, amsthm, graphicx, url, enumerate; the wrapper is the lane's pinned \texttt{amsart} editorial wrapper at 10pt on A4, which restates the theorem-like environments the retained text needs and adds the article's own macro definitions that the retained text needs (\texttt{\textbackslash CZ}, \texttt{\textbackslash mI}, \texttt{\textbackslash zdu}), with extra wrapper packages (bm, bbm, dsfont, xspace, cancel, mathtools). The delivered numbering is the unit's own. Assets: no retained span contains a figure environment, a graphics-inclusion command, a PDF-page inclusion, a table or a TikZ picture; no image file is copied into this package, the package declares no asset dependency and no image file is redistributed with it. Licence: version arXiv:2502.05678v1 of this article is distributed under the Creative Commons Attribution 4.0 International licence, as recorded in the arXiv licence metadata for this exact version and shown on the abstract page https://arxiv.org/abs/2502.05678v1. A full rights notice is delivered at licenses/arxiv-2502.05678-CC-BY-4.0.txt. Characters: every retained excerpt is pure ASCII, so the delivered engine, XeLaTeX with its default Latin Modern text font, reports no missing character.
\begin{theorem}\label{laplacian eigenvalue}
Let $G=(V,E)$ be a simple finite graph on $m$ vertices $\{v_1, \ldots, v_m\}$ and suppose that for some positive integer $d$, the graph has a unique vertex $v_k$ of degree $d$.  It follows that the zeon combinatorial Laplacian $\Lambda$ of $G$ has a spectrally simple eigenvalue $\lambda$ satisfying \begin{equation*}
\lambda = d + \sum_{I\subseteq [m]}\omega_I\zeta_{I}\prod_{j\in I\setminus\{k\}}(d-\deg(v_j))^{-1},
\end{equation*}
where $(-1)^{|I|}\omega_I$ is the number of cycles based at $v_k$ on vertices indexed by $I$.
\end{theorem}

\begin{proof}
Begin by relabeling the vertices $\{v_1, \ldots, v_m\}$ if necessary so that $v_m$ is the unique vertex of degree $d$, and let $\Lambda$ be the zeon combinatorial Laplacian of $G$ after relabeling.  Since $v_m$ is the unique vertex of degree $d$, the first $m-1$ main diagonal elements of $\lambda\mI-\Lambda$ are invertible, while the last main diagonal element is nilpotent.   In particular, letting $d_j$ denote the degree of vertex $v_j$ and writing $\ell_j=\lambda-d_j$ for $j=1, \ldots, m-1$  , we have
\begin{equation}\label{form to triangularize}
\lambda\mI-\Lambda=\left(\begin{matrix}
\ell_1&\varepsilon_{12}\zeta_{\{2\}}&\cdots&\cdots&\varepsilon_{1m}\zeta_{\{m\}}\\
\varepsilon_{12}\zeta_{\{1\}}&\ell_2&\varepsilon_{13}\zeta_{\{3\}}&\cdots&\varepsilon_{2m}\zeta_{\{m\}}\\
\varepsilon_{13}\zeta_{\{1\}}&\varepsilon_{23}\zeta_{\{2\}}&\ell_3&\cdots&\varepsilon_{3m}\zeta_{\{m\}}\\
\vdots&\vdots&\vdots&\ddots&\vdots\\
\varepsilon_{1m}\zeta_{\{1\}}&\cdots&\cdots&\cdots&\zdu \lambda
\end{matrix}\right),  
 \end{equation}
where $\varepsilon_{ij}$ denotes the function $V\times V\to \{0,1\}$ defined by 
\begin{equation*}
\varepsilon_{ij}=\begin{cases}
1&\text{\rm if }\{v_i,v_j\}\in E,\\
0&\text{\rm otherwise.}
\end{cases}
\end{equation*}


It follows that by simply adding zeon multiples of rows to each other, $\lambda\mI-\Lambda$ can be placed into upper triangular form  \begin{equation}\label{triangular form}
\tau(\lambda\mI-\Lambda)=\left(\begin{matrix}
\alpha_1&*&\cdots&*&\eta_1\\
0&\alpha_2&\cdots&*&\eta_2\\
0&0&\ddots&\vdots&\vdots\\
\vdots&\vdots&0&\alpha_{m-1}&\eta_{m-1}\\
0&0&\cdots&0&\eta_m
\end{matrix}\right),  
 \end{equation}
where $\alpha_\ell\in\CZ^\times$ and $\eta_\ell\in\CZ^\circ$ for $\ell=1, \ldots, m$.  Further, since the only invertible elements of $\lambda\mI-\Lambda$ lie along the main diagonal, the only nilpotent column of $\lambda\mI-\Lambda$ is the $m$th column.   Since zeros are obtained below each element along the main diagonal by adding a zeon multiple of one row to another, the determinant of the resulting matrix is unchanged.
Hence, 
\begin{equation*}
0=\vert \lambda\mI-\Lambda\vert=\vert \tau(\lambda\mI-\Lambda)\vert=\eta_m\prod_{\ell=1}^{m-1}\alpha_\ell,
\end{equation*}
implying $\eta_m=0$.


Now set $\tau_0=\lambda\mI-\Lambda$ and let $\tau_k$ be the matrix obtained after $k$ steps of the triangularization process for $k=1, \ldots, m-1$.  

\begin{description}
\item [{\bf Claim 1.}] The off-diagonal elements of $\tau_k$ are given by the following:
\begin{equation}\label{off-diag}
(\tau_k)_{ij}=\begin{cases}
0&  i>j, j\le k;\\
\varepsilon_{ij}\zeta_{\{j\}}-\sum_{I\subseteq [k]}\gamma_{I\cup\{j\}}\zeta_{I\cup\{j\}}\prod_{t\in I\cap[k]}{\ell_t}^{-1}& i > j > k;\\
\varepsilon_{ij}\zeta_{\{j\}}-\sum_{I\subseteq [k-1]}\gamma_{I\cup\{j\}}\zeta_{I\cup\{j\}}\prod_{t\in I\cap[i-1]}{\ell_t}^{-1}&  i< j, j>k;\\
\varepsilon_{ij}\zeta_{\{j\}}-\sum_{I\subseteq [i-1]}\gamma_{I\cup\{j\}}\zeta_{I\cup\{j\}}\prod_{t\in I\cap[i-1]}{\ell_t}^{-1} &i<j\le k;
\end{cases}
\end{equation}where $(-1)^{|I\cup\{j\}|}\gamma_{I\cup\{j\}}$ is the number of paths from $v_i$ to $v_j$ on vertices indexed by $I\cup\{j\}$ (excluding the initial vertex $v_i$; i.e.,  $i\notin I\cup\{j\}$).  
\end{description}

Proof of the claim is by induction on $k$.  When $k=0$, the result holds by definition of $\tau_0$.  Now suppose \eqref{off-diag} is valid for all off-diagonal entries of $(\tau_{k-1})$.   
We now proceed by cases.

\begin{description}
\item [{\bf Case 1. ($i>j$ and $j\le k$)}]
Zeros are obtained below $(\tau_{k-1})_{kk}$ via Gaussian elimination using the inverse $((\tau_{k-1})_{kk})^{-1}$ in the usual way since we assume all diagonal entries (except possibly the last one) are invertible.  This gives $(\tau_k)_{ij}=0$ whenever $j<i$ and $j\le k$.
\end{description}

\begin{description}
\item [{\bf Case 2. ($i>j> k$)}]

We consider now those entries below the main diagonal that have not yet been zeroed out.  If $v_i$ is not adjacent to $v_k$, then $(\tau_{k-1})_{ik}=0$, so that $(\tau_k)_{ij}=(\tau_{k-1})_{ij}$; i.e., no change is made to row $i$ in the triangularization $\tau_{k-1}\to\tau_k$.  On the other hand, if $v_i$ is adjacent to $v_k$, then $\varepsilon_{ik}\ne 0$ and \begin{equation}
(\tau_k)_{ij}=(\tau_{k-1})_{ij}-\frac{(\tau_{k-1})_{ik}}{(\tau_{k-1})_{kk}}(\tau_{k-1})_{kj}.
\end{equation}  
Writing $(\tau_{k-1})_{kk}=\ell_k-\xi$, where $\xi\in\CZ^\circ$ is a $\CZ$-linear combination of elements annihilated by $\zeta_{\{k\}}$, it follows that $((\tau_{k-1})_{kk})^{-1}=\frac{1}{\ell_k}+\frac{1}{{\ell_k}^2}\xi$.  
Further, by the inductive hypothesis, $(\tau_{k-1})_{ik}$ is a $\CZ$-linear combination of elements annihilated by $\zeta_{\{k\}}$. Consequently, \begin{equation*}
\frac{(\tau_{k-1})_{ik}}{(\tau_{k-1})_{kk}}=(\tau_{k-1})_{ik}\left(\frac{1}{\ell_k}+\frac{1}{{\ell_k}^2}\xi\right)=\frac{(\tau_{k-1})_{ik}}{\ell_k}.
\end{equation*}
Hence, \begin{align*}
(\tau_k)_{ij}&=(\tau_{k-1})_{ij}-\frac{(\tau_{k-1})_{ik}}{\ell_k}(\tau_{k-1})_{kj}\\
&=\varepsilon_{ij}\zeta_{\{j\}}-\sum_{I\subseteq [k-1]}\gamma_{I\cup\{j\}}\zeta_{I\cup\{j\}}\prod_{t\in I\cap[k-1]}{\ell_t}^{-1}\\
&\hskip10pt-{\ell_k}^{-1}(\tau_{k-1})_{ik}(\tau_{k-1})_{kj}\\
&=\varepsilon_{ij}\zeta_{\{j\}}-\sum_{I\subseteq [k-1]}\gamma_{I\cup\{j\}}\zeta_{I\cup\{j\}}\prod_{t\in I\cap[k-1]}{\ell_t}^{-1}\\
&\hskip10pt-{\ell_k}^{-1}\left(\varepsilon_{ik}\zeta_{\{k\}}-\sum_{I\subseteq [k-1]}{\gamma_{I\cup\{k\}}}'\zeta_{I\cup\{k\}}\prod_{t\in I\cap[k-1]}{\ell_t}^{-1}\right)\\
&\hskip10pt\times\left(\varepsilon_{kj}\zeta_{\{j\}}-\sum_{I\subseteq [k-1]}{\gamma_{I\cup\{j\}}}''\zeta_{I\cup\{j\}}\prod_{t\in I\cap[k-1]}{\ell_t}^{-1}\right),
\end{align*}
where $(-1)^{|I\cup\{j\}|}\gamma_{I\cup\{j\}}$ is the number of paths from $v_i$ to $v_j$ on vertices indexed by $I\cup\{j\}$, $(-1)^{|I|\cup\{k\}}{\gamma_{I\cup\{k\}}}'$ is the number of paths from $v_i$ to $v_k$ on vertices indexed by $I\cup\{k\}$, and $(-1)^{|I|\cup\{j\}}{\gamma_{I\cup\{j\}}}''$ is the number of paths from $v_k$ to $v_j$ on vertices indexed by $I\cup\{j\}$, respectively.

Expanding the right-hand side of the final equality, we obtain 
\begin{align*}
(\tau_k)_{ij}&=\varepsilon_{ij}\zeta_{\{j\}}-\sum_{I\subseteq [k-1]}\gamma_{I\cup\{j\}}\zeta_{I\cup\{j\}}\prod_{t\in I\cap[k-1]}{\ell_t}^{-1}\\
&\hskip10pt-{\ell_k}^{-1}\varepsilon_{ik}\varepsilon_{kj}\zeta_{\{j,k\}}\\
&\hskip10pt+{\ell_k}^{-1}\varepsilon_{ik}\zeta_{\{k\}}\sum_{I\subseteq [k-1]}{\gamma_{I\cup\{j\}}}''\zeta_{I\cup\{j\}}\prod_{t\in I\cap[k-1]}{\ell_t}^{-1}\\
&\hskip10pt+{\ell_k}^{-1}\varepsilon_{kj}\zeta_{\{j\}}\sum_{I\subseteq [k-1]}{\gamma_{I\cup\{k\}}}'\zeta_{I\cup\{k\}}\prod_{t\in I\cap[k-1]}{\ell_t}^{-1}\\
&\hskip10pt-{\ell_k}^{-1}\left[\sum_{I\subseteq [k-1]}{\gamma_{I\cup\{k\}}}'\zeta_{I\cup\{k\}}\prod_{t\in I\cap[k-1]}{\ell_t}^{-1}\right]\\
&\hskip10pt\times\left[\sum_{I\subseteq [k-1]}{\gamma_{I\cup\{j\}}}''\zeta_{I\cup\{j\}}\prod_{t\in I\cap[k-1]}{\ell_t}^{-1}\right].
\end{align*}

Note that $\varepsilon_{ij}\zeta_{\{j\}}$ represents the unique one-step path from $v_i$ to $v_j$ when $v_i$ and $v_j$ are adjacent.  Further, $-{\ell_k}^{-1}\varepsilon_{ik}\varepsilon_{kj}\zeta_{\{j,k\}}$ represents the unique two-step path $v_i\to v_k\to v_j$ when $\{v_i,v_k\}, \{v_k,v_j\}\in E$.  

The quantity \begin{equation*}
+{\ell_k}^{-1}\varepsilon_{ik}\zeta_{\{k\}}\sum_{I\subseteq [k-1]}{\gamma_{I\cup\{j\}}}''\zeta_{I\cup\{j\}}\prod_{t\in I\cap[k-1]}{\ell_t}^{-1}
\end{equation*}
represents paths $v_i\to v_j$ on vertices indexed by subsets of $[k]$ passing through vertex $v_k$ in the second step.

The quantity \begin{equation*}
{\ell_k}^{-1}\varepsilon_{kj}\zeta_{\{j\}}\sum_{I\subseteq [k-1]}{\gamma_{I\cup\{k\}}}'\zeta_{I\cup\{k\}}\prod_{t\in I\cap[k-1]}{\ell_t}^{-1}
\end{equation*}
represents paths $v_i\to v_j$ on vertices indexed by subsets of $[k]$ passing through vertex $v_k$ in the penultimate step.

Finally, the quantity 
\begin{align*}
-{\ell_k}^{-1}\left[\sum_{I\subseteq [k-1]}{\gamma_{I\cup\{k\}}}'\zeta_{I\cup\{k\}}\prod_{t\in I\cap[k-1]}{\ell_t}^{-1}\right]\\
 \times
\left[\sum_{I\subseteq [k-1]}{\gamma_{I\cup\{j\}}}''\zeta_{I\cup\{j\}}\prod_{t\in I\cap[k-1]}{\ell_t}^{-1}\right]
\end{align*}
represents the product of walks $v_i\to v_k$ with walks from $v_k\to v_j$ on vertices indexed by a subset of $[k]$ which pass though $v_k$ in some intermediate step.  



\end{description}

\begin{description}
\item [{\bf Case 3. ($i<j$ and $j>k$)}]

Considering entries $(\tau_k)_{ij}$ above the main diagonal with $j>k$,  the proof proceeds in similar fashion to that of Case 2.   In particular, 
 \begin{align*}
(\tau_k)_{ij}&=(\tau_{k-1})_{ij}-\frac{(\tau_{k-1})_{ik}}{\ell_k}(\tau_{k-1})_{kj}\\
&=\varepsilon_{ij}\zeta_{\{j\}}-\sum_{I\subseteq [k-2]}\gamma_{I\cup\{j\}}\zeta_{I\cup\{j\}}\prod_{t\in I\cap[i-1]}{\ell_t}^{-1}\\
&\hskip10pt-{\ell_k}^{-1}(\tau_{k-1})_{ik}(\tau_{k-1})_{kj}\\
&=\varepsilon_{ij}\zeta_{\{j\}}-\sum_{I\subseteq [k-2]}\gamma_{I\cup\{j\}}\zeta_{I\cup\{j\}}\prod_{t\in I\cap[i-1]}{\ell_t}^{-1}\\
&\hskip10pt-{\ell_k}^{-1}\left(\varepsilon_{ik}\zeta_{\{k\}}-\sum_{I\subseteq [i-1]}{\gamma_{I\cup\{k\}}}'\zeta_{I\cup\{k\}}\prod_{t\in I\cap[i-1]}{\ell_t}^{-1}\right)\\
&\hskip10pt\times\left(\varepsilon_{kj}\zeta_{\{j\}}-\sum_{I\subseteq [k-2]}{\gamma_{I\cup\{j\}}}''\zeta_{I\cup\{j\}}\prod_{t\in I\cap[k-1]}{\ell_t}^{-1}\right),
\end{align*}
where $(-1)^{|I\cup\{j\}|}\gamma_{I\cup\{j\}}$ is the number of paths and PWICs from $v_i$ to $v_j$ on vertices indexed by $I\cup\{j\}$, $(-1)^{|I|\cup\{k\}}{\gamma_{I\cup\{k\}}}'$ is the number of paths and PWICs from $v_i$ to $v_k$ on vertices indexed by $I\cup\{k\}$, and $(-1)^{|I|\cup\{j\}}{\gamma_{I\cup\{j\}}}''$ is the number of paths and PWICs from $v_k$ to $v_j$ on vertices indexed by $I\cup\{j\}$, respectively.
Expanding the right-hand side of the final equality, we obtain 
\begin{align*}
(\tau_k)_{ij}&=\varepsilon_{ij}\zeta_{\{j\}}-\sum_{I\subseteq [k-2]}\gamma_{I\cup\{j\}}\zeta_{I\cup\{j\}}\prod_{t\in I\cap[i-1]}{\ell_t}^{-1}\\
&\hskip10pt-{\ell_k}^{-1}\varepsilon_{ik}\varepsilon_{kj}\zeta_{\{j,k\}}\\
&\hskip10pt+{\ell_k}^{-1}\varepsilon_{ik}\zeta_{\{k\}}\sum_{I\subseteq [k-2]}{\gamma_{I\cup\{j\}}}''\zeta_{I\cup\{j\}}\prod_{t\in I\cap[k-1]}{\ell_t}^{-1}\\
&\hskip10pt+{\ell_k}^{-1}\varepsilon_{kj}\zeta_{\{j\}}\sum_{I\subseteq [i-1]}{\gamma_{I\cup\{k\}}}'\zeta_{I\cup\{k\}}\prod_{t\in I\cap[i-1]}{\ell_t}^{-1}\\
&\hskip10pt-{\ell_k}^{-1}\left[\sum_{I\subseteq [k-2]}{\gamma_{I\cup\{k\}}}'\zeta_{I\cup\{k\}}\prod_{t\in I\cap[k-1]}{\ell_t}^{-1}\right]\\
&\hskip10pt\times\left[\sum_{I\subseteq [i-1]}{\gamma_{I\cup\{j\}}}''\zeta_{I\cup\{j\}}\prod_{t\in I\cap[i-1]}{\ell_t}^{-1}\right].
\end{align*}


Recall that $\varepsilon_{ij}\zeta_{\{j\}}$ represents the unique one-step path from $v_i$ to $v_j$ when $v_i$ and $v_j$ are adjacent.  Further, $-{\ell_k}^{-1}\varepsilon_{ik}\varepsilon_{kj}\zeta_{\{j,k\}}$ represents the unique two-step path $v_i\to v_k\to v_j$ when $\{v_i,v_k\}, \{v_k,v_j\}\in E$.  

Since $(-1)^{|I|\cup\{j\}}{\gamma_{I\cup\{j\}}}''$ is the number of paths and PWICs $v_k\to v_j$, and $\varepsilon_{ik}=1$ if and only if $v_i$ is adjacent to $v_k$, the quantity \begin{equation*}
{\ell_k}^{-1}\varepsilon_{ik}\zeta_{\{k\}}\sum_{I\subseteq [k-2]}{\gamma_{I\cup\{j\}}}''\zeta_{I\cup\{j\}}\prod_{t\in I\cap[k-1]}{\ell_t}^{-1}
\end{equation*}
represents paths and PWICs $v_i\to v_j$ on vertices indexed by subsets of $[k]$ passing through vertex $v_k$ in the second step.

By similar reasoning, the quantity \begin{equation*}
{\ell_k}^{-1}\varepsilon_{kj}\zeta_{\{j\}}\sum_{I\subseteq [i-1]}{\gamma_{I\cup\{k\}}}'\zeta_{I\cup\{k\}}\prod_{t\in I\cap[i-1]}{\ell_t}^{-1}\end{equation*}
represents paths $v_i\to v_j$ on vertices indexed by subsets of $[k]$ passing through vertex $v_k$ in the penultimate step.

Finally, the quantity 
\begin{align*}
-{\ell_k}^{-1}\left[\sum_{I\subseteq [k-2]}{\gamma_{I\cup\{k\}}}'\zeta_{I\cup\{k\}}\prod_{t\in I\cap[k-1]}{\ell_t}^{-1}\right]\\
\times\left[\sum_{I\subseteq [i-1]}{\gamma_{I\cup\{j\}}}''\zeta_{I\cup\{j\}}\prod_{t\in I\cap[i-1]}{\ell_t}^{-1}\right]
\end{align*}
represents the product of walks $v_i\to v_k$ with walks from $v_k\to v_j$ on vertices indexed by a subset of $[k]$ which pass though $v_k$ in some intermediate step.  Summing the quantities establishes the stated result for Case 3.
\end{description}

\begin{description}
\item [{\bf Case 4. ($i<j\le k$)}]
Case 4 follows from Case 3 by observing that once zeros are obtained below diagonal $(\tau_k)_{kk}$, rows 1 through $k$ remain fixed; i.e., $(\tau_n)_{ij}=(\tau_k)_{ij}$ for $1\le i\le k$ and $k\le n\le m$.   Hence, the only values ${\ell_t}^{-1}$ appearing in terms of $(\tau_k)_{ij}$ are indexed by $t\in [i-1]$. 
\end{description}

This concludes the proof of Claim 1.  We now turn to elements along the main diagonal of $\tau_k$.

\begin{description}
\item [{\bf Claim 2.}] The diagonal elements of $\tau_k$ are given by the following:
\begin{equation}\label{diag elts}
(\tau_k)_{jj}=\begin{cases}\ell_j-\sum_{I\subseteq [k]}\omega_{I\cup\{j\}}\zeta_{I\cup\{j\}}\prod_{t\in I\cap[k]}{\ell_t}^{-1}&j>k,\\
\ell_j-\sum_{I\subseteq [j]}{\omega_I}'\zeta_I\prod_{t\in I\cap[j-1]}{\ell_t}^{-1}&j\le k,
\end{cases}
\end{equation} where $(-1)^{|I\cup\{j\}|}\omega_{I\cup\{j\}}$ is the number of cycles based at $v_j$ on vertices indexed by $I\cup\{j\}$ and $(-1)^{|I|}{\omega_I}'$ is the number of cycles based at $v_j$ on vertices indexed by $I$.  
\end{description}

Proof of this claim is again by induction on $k$.  When $k=0$, the result holds by definition of $\tau_0$.  Now suppose \eqref{off-diag} is valid for all off-diagonal entries of $(t_{k-1})$ and that \eqref{diag elts} is valid for diagonal elements of $(\tau_{k-1})$.   If $v_k$ is not adjacent to $v_j$, we see that $(\tau_k)_{jj}=(\tau_{k-1})_{jj}$, and there is nothing to prove.  On the other hand, if $v_k$ is adjacent to $v_j$, we have 
\begin{equation}\label{diagonal recurrence}
(\tau_k)_{jj}=(\tau_{k-1})_{jj}-\frac{(\tau_{k-1})_{jk}}{(\tau_{k-1})_{kk}}(\tau_{k-1})_{kj}.
\end{equation}  

We now continue by cases.

\begin{description}
\item [{\bf Case 1. ($j>k$)}]
When $j>k$, \eqref{diagonal recurrence} becomes
\begin{align*}
(\tau_k)_{jj}&=\ell_j-\sum_{I\subseteq [k-1]\cup\{j\}}\omega_I\zeta_I\prod_{t\in I}{\ell_t}^{-1}-\frac{1}{\ell_k}(\tau_{k-1})_{jk}(\tau_{k-1})_{kj}
\end{align*}
where $(-1)^{|I|}\omega_I$ is the number of cycles based at $v_j$ on vertices indexed by $I\subseteq[k-1]\cup\{j\}$ and \eqref{off-diag} gives
\begin{align*}
(\tau_{k-1})_{jk}(\tau_{k-1})_{kj}&=\left(\zeta_{\{k\}}-\sum_{I\subseteq [k-1]\cup\{k\}}\gamma_I\zeta_I\prod_{t\in I\setminus\{k-1\}}{\ell_t}^{-1}\right)\\
&\hskip10pt\times\left(\zeta_{\{j\}}-\sum_{I\subseteq [k-1]\cup\{j\}}{\gamma_I}'\zeta_I\prod_{t\in I\setminus\{k-1\}}{\ell_t}^{-1}\right)\\
&=\zeta_{\{j,k\}}-\zeta_{\{k\}}\sum_{I\subseteq [k-1]\cup\{j\}}{\gamma_I}'\zeta_I\prod_{t\in I\setminus\{k-1\}}{\ell_t}^{-1}\\
&\hskip10pt -\zeta_{\{j\}}\sum_{I\subseteq [k-1]\cup\{k\}}\gamma_I\zeta_I\prod_{t\in I\setminus\{k-1\}}{\ell_t}^{-1}\\
&\hskip10pt+\left(\sum_{I\subseteq [k-1]\cup\{k\}}\gamma_I\zeta_I\prod_{t\in I\setminus\{k-1\}}{\ell_t}^{-1}\right)\\
&\hskip10pt\times\left(\sum_{I\subseteq [k-1]\cup\{j\}}{\gamma_I}'\zeta_I\prod_{t\in I\setminus\{k-1\}}{\ell_t}^{-1}\right),
\end{align*}
where $(-1)^{|I|}\gamma_I$ denotes the number of paths from $v_j$ to $v_k$ on vertices indexed by $I\subseteq [k]$ and $(-1)^{|I|}{\gamma_I}'$ denotes the number of paths from $v_k$ to $v_j$ on vertices indexed by $I\subseteq [k-1]\cup \{j\}$.  Now $\zeta_{\{j,k\}}$ represents a 2-cycle on vertices $\{v_j, v_k\}$.  The quantity \begin{equation*}
\zeta_{\{k\}}\sum_{I\subseteq [k-1]\cup\{j\}}{\gamma_I}'\zeta_I\prod_{t\in I\setminus\{k-1\}}{\ell_t}^{-1}=\sum_{I\subseteq [k]\cup\{j\}}{\gamma_I}'\zeta_I\prod_{t\in I\setminus\{k-1\}}{\ell_t}^{-1}
\end{equation*}
represents cycles based at $v_j$ whose first step is to $v_j\to v_k$.  Similarly, 
\begin{equation*}
\zeta_{\{j\}}\sum_{I\subseteq [k-1]\cup\{k\}}\gamma_I\zeta_I\prod_{t\in I\setminus\{k-1\}}{\ell_t}^{-1}
\end{equation*}
represents cycles based at $v_j$ whose last step is $v_k\to v_j$.
Finally, fixing $I\subseteq[k]$ and $J\subseteq[k-1]\cup\{j\}$,  the nonzero terms of the product of last terms are of the form
\begin{align}
\left[\gamma_I\zeta_I\prod_{t\in I\setminus\{k-1\}}{\ell_t}^{-1}\right]
\left[{\gamma_J}'\zeta_J\prod_{s\in J\setminus\{k-1\}}{\ell_s}^{-1}\right]\nonumber\\
=\gamma_I{\gamma_J}'\zeta_{I\cup J}\prod_{r\in (I\cup J)\setminus\{k-1\}}{\ell_r}^{-1},\end{align}
provided $I\cap J=\varnothing$.  This represents all remaining cycles $v_j\to v_k\to v_j$ on vertices indexed by $I\cup J$.  Now $(-1)^{|I\cup J|}\gamma_I{\gamma_J}'$ denotes the number of cycles based at $v_j$ on vertices indexed by $I\cup J$ including $v_k$ at an intermediate step.  
\end{description}

\begin{description}
\item [{\bf Case 2. ($j\le k$)}]
Recalling that $(\tau_\ell)_{ij}=(\tau_k)_{ij}$ for $1\le \ell\le k$, the only values ${\ell_t}^{-1}$ appearing in terms of $(\tau_k)_{jj}$ are indexed by $t\in [j-1]$.  Otherwise, the proof proceeds as the proof of Case 1.
\end{description}
This concludes the proof of Claim 2.  Letting $k=m-1$ and $j=m$ in \eqref{diag elts}, we now see that  \begin{equation}\label{last diagonal}
(\tau_{m-1})_{mm}=\ell_m-\sum_{I\subseteq [m]}\omega_I\zeta_I\prod_{j\in I\setminus\{m\}}{\ell_j}^{-1},
\end{equation}
where $(-1)^{|I|}\omega_I$ is the number of cycles based at $v_m$ on vertices indexed by $I\subseteq [m]$.  In particular, $m\in I$ for each nonzero term of \eqref{last diagonal}.

Referring to \eqref{triangular form}, each $\eta_j$ is a $\CZ$-linear combination of elements annihilated by $\zeta_{\{m\}}$.  Further, $\eta_m=0$ implies $\zdu\lambda$ is a $\CZ$-linear combination of $\eta_j$s, so $\zeta_{\{m\}}$ annihilates $\zdu\lambda$.  Consequently,  $(\zdu \lambda)^2=0$, and $\eta_j\zdu\lambda=0$ for $j=1, \ldots, m-1$.  Moreover,  it follows that \begin{equation*}
{\ell_j}^{-1}=\frac{1}{d-d_j}-\frac{1}{(d-d_j)^2}\zdu\lambda,
\end{equation*}
and that $\zeta_I\zdu\lambda =0$ when $m\in I$. 

The proof of the theorem is now completed by recalling that $\ell_m=\zdu \lambda$ and $(\tau_{m-1})_{mm}=0$.   In particular, 
\begin{align*}
\zdu\lambda&=\sum_{I\subseteq [m]}\omega_I\zeta_I\prod_{j\in I\setminus\{m\}}{\ell_j}^{-1}\\
&=\sum_{I\subseteq [m]}\omega_I\zeta_{I}\prod_{j\in I\setminus\{m\}}\left(\frac{1}{d-d_j}-\frac{1}{(d-d_j)^2}\zdu\lambda\right)\\
&=\sum_{I\subseteq [m]}\omega_I\zeta_{I}\prod_{j\in I\setminus\{m\}}\frac{1}{d-d_j}\\
&=\lambda-d,
\end{align*}
where $(-1)^{|I|}\omega_I$ is the number of cycles based at $v_m$ on vertices indexed by $I\subseteq [m]$.  

\end{proof}

\end{document}
